\documentclass[12pt]{article}
\usepackage{amsmath, amssymb, amsthm, geometry, hyperref}
\geometry{a4paper, margin=1in}

\title{Linear Algebra Problem Sets: Sections 1.3 \& 1.4}
\author{Curated Practice \& Solutions}
\date{}

\begin{document}

\maketitle
\tableofcontents
\newpage

% ==========================================
% SECTION: CONCEPT CHECKLIST
% ==========================================
\section{Concept Checklist \& Question Mapping}

\subsection{Section 1.3: Vector Equations}
\begin{itemize}
    \item \textbf{Vector Arithmetic \& Linear Combinations:} Computing $u+v$, $cu-dv$, etc. \textit{(Questions: 1, 2, 3)}
    \item \textbf{Translation (System $\leftrightarrow$ Vector Equation):} Converting between systems of equations and vector representations. \textit{(Questions: 4, 5, 6, 7, 8, 9)}
    \item \textbf{Membership in a Span:} Determining if a vector $b$ is a linear combination of given vectors (setting up augmented matrix and checking consistency). \textit{(Questions: 10, 11, 12, 13, 14, 15)}
    \item \textbf{Geometric Interpretation:} Understanding spans as points, lines, or planes, and identifying collinearity. \textit{(Questions: 16, 17, 18)}
    \item \textbf{Set vs. Span:} Distinguishing between a vector belonging to a finite set versus belonging to the infinite span of that set. \textit{(Questions: 19, 20)}
    \item \textbf{Parameterized Span:} Finding an unknown variable $h$ such that $b$ is in the span. \textit{(Questions: 21, 22, 23, 24)}
    \item \textbf{Real-World Applications:} Center of mass, resource/diet optimization. \textit{(Questions: 25, 26, 27)}
    \item \textbf{Spanning all of $\mathbb{R}^n$:} Determining if a set of vectors spans the entire real space. \textit{(Questions: 28, 29)}
    \item \textbf{Advanced Diagnostic ("Find the Flaw"):} Identifying logical/algebraic errors. \textit{(Question: 30)}
\end{itemize}

\subsection{Section 1.4: The Matrix Equation $Ax=b$}
\begin{itemize}
    \item \textbf{Computing $Ax$ (Definition \& Row-Vector Rule):} Multiplying a matrix by a vector using both linear combinations of columns and the dot-product rule. \textit{(Questions: 1, 2, 3, 4, 5, 6)}
    \item \textbf{Dimensional Compatibility:} Checking if a matrix-vector product is defined. \textit{(Questions: 7, 8)}
    \item \textbf{Translation (System/Vector $\leftrightarrow$ Matrix Equation):} Converting systems and vector equations into $Ax=b$. \textit{(Questions: 9, 10, 11, 12, 13, 14)}
    \item \textbf{Solving $Ax=b$:} Row reducing augmented matrices to extract vector solutions (unique, infinite, or no solution). \textit{(Questions: 15, 16, 17, 18, 19)}
    \item \textbf{Existence for all $b$ (Theorem 4):} Determining if $A$ spans $\mathbb{R}^m$ by looking for pivots in every row. \textit{(Questions: 20, 21, 22)}
    \item \textbf{Linearity Properties:} Using $A(u+v) = Au + Av$ and $A(cu) = cAu$. \textit{(Questions: 23, 24, 25)}
    \item \textbf{Parameterized Consistency:} Finding conditions on $h$ or $k$ for the system to be consistent. \textit{(Questions: 26, 27)}
    \item \textbf{Real-World Applications:} Derivative replication in finance. \textit{(Question: 28)}
    \item \textbf{Generalization \& Diagnostics:} Proving structural limits (e.g., tall matrices) and finding flaws. \textit{(Questions: 29, 30)}
\end{itemize}

\newpage

% ==========================================
% SECTION: PROBLEM SETS
% ==========================================
\section{Problem Sets}

\subsection{Homework 1.3: Vector Equations}

\textbf{Vector Arithmetic \& Linear Combinations} \\
Let $u = \begin{bmatrix} 2 \\ -1 \end{bmatrix}$, $v = \begin{bmatrix} 3 \\ 4 \end{bmatrix}$, and $w = \begin{bmatrix} -1 \\ 2 \end{bmatrix}$. 
\begin{enumerate}
    \item Compute $u + 2v$.
    \item Compute $3u - v + 2w$.
    \item Find scalars $c_1, c_2$ such that $c_1 u + c_2 w = \begin{bmatrix} 7 \\ -5 \end{bmatrix}$.
\end{enumerate}

\textbf{Translation (System $\leftrightarrow$ Vector Equation)} \\
Write a vector equation that is equivalent to the given system of equations:
\begin{enumerate}
\setcounter{enumi}{3}
    \item $3x_1 - x_2 = 5$ \\ $x_1 + 4x_2 = -2$
    \item $x_1 + 2x_2 - x_3 = 0$ \\ $2x_1 - x_2 + 3x_3 = 4$ \\ $-x_1 + 5x_3 = 7$
    \item $4x_1 - 2x_2 = 1$ \\ $x_1 + x_2 = 3$ \\ $-2x_1 + 5x_2 = 0$
\end{enumerate}

Write a system of linear equations equivalent to the given vector equation:
\begin{enumerate}
\setcounter{enumi}{6}
    \item $x_1 \begin{bmatrix} 1 \\ -3 \end{bmatrix} + x_2 \begin{bmatrix} 4 \\ 2 \end{bmatrix} = \begin{bmatrix} 6 \\ -1 \end{bmatrix}$
    \item $x_1 \begin{bmatrix} 2 \\ 0 \\ 1 \end{bmatrix} + x_2 \begin{bmatrix} -1 \\ 5 \\ 3 \end{bmatrix} = \begin{bmatrix} 0 \\ 2 \\ -4 \end{bmatrix}$
    \item $x_1 \begin{bmatrix} 1 \\ 1 \\ -2 \end{bmatrix} + x_2 \begin{bmatrix} 0 \\ -3 \\ 1 \end{bmatrix} + x_3 \begin{bmatrix} 4 \\ 2 \\ 0 \end{bmatrix} = \begin{bmatrix} 8 \\ 9 \\ 10 \end{bmatrix}$
\end{enumerate}

\textbf{Membership in a Span} \\
Determine if vector $b$ is a linear combination of the columns of the given matrix / set of vectors.
\begin{enumerate}
\setcounter{enumi}{9}
    \item $a_1 = \begin{bmatrix} 1 \\ 2 \end{bmatrix}, a_2 = \begin{bmatrix} -3 \\ 4 \end{bmatrix}, b = \begin{bmatrix} -5 \\ 10 \end{bmatrix}$
    \item $a_1 = \begin{bmatrix} 2 \\ 4 \end{bmatrix}, a_2 = \begin{bmatrix} -1 \\ -2 \end{bmatrix}, b = \begin{bmatrix} 3 \\ 5 \end{bmatrix}$
    \item $a_1 = \begin{bmatrix} 1 \\ 0 \\ -1 \end{bmatrix}, a_2 = \begin{bmatrix} 2 \\ 1 \\ 3 \end{bmatrix}, b = \begin{bmatrix} 4 \\ 1 \\ 1 \end{bmatrix}$
    \item $a_1 = \begin{bmatrix} 1 \\ -1 \\ 2 \end{bmatrix}, a_2 = \begin{bmatrix} 3 \\ 0 \\ 1 \end{bmatrix}, b = \begin{bmatrix} 5 \\ -2 \\ 6 \end{bmatrix}$
    \item $a_1 = \begin{bmatrix} 1 \\ 0 \\ 0 \end{bmatrix}, a_2 = \begin{bmatrix} 1 \\ 1 \\ 0 \end{bmatrix}, a_3 = \begin{bmatrix} 1 \\ 1 \\ 1 \end{bmatrix}, b = \begin{bmatrix} 4 \\ -2 \\ 7 \end{bmatrix}$
    \item $a_1 = \begin{bmatrix} 1 \\ 2 \\ 3 \end{bmatrix}, a_2 = \begin{bmatrix} -1 \\ -2 \\ -3 \end{bmatrix}, a_3 = \begin{bmatrix} 2 \\ 4 \\ 6 \end{bmatrix}, b = \begin{bmatrix} 1 \\ 1 \\ 1 \end{bmatrix}$
\end{enumerate}

\textbf{Geometric Interpretation} 
\begin{enumerate}
\setcounter{enumi}{15}
    \item Describe the geometric span of $v_1 = \begin{bmatrix} 3 \\ -6 \\ 9 \end{bmatrix}$ and $v_2 = \begin{bmatrix} -1 \\ 2 \\ -3 \end{bmatrix}$.
    \item Let $v_1, v_2, v_3$ be three non-zero, mutually non-collinear vectors in $\mathbb{R}^3$ where $v_3 = v_1 + v_2$. Does their span form a line, a plane, or all of $\mathbb{R}^3$?
    \item If a set of vectors includes the zero vector $0$, does the zero vector change the geometric span of the other vectors? Explain.
\end{enumerate}

\textbf{Set vs. Span} \\
Let $v_1 = \begin{bmatrix} 1 \\ 3 \end{bmatrix}, v_2 = \begin{bmatrix} 2 \\ -1 \end{bmatrix}$, and $b = \begin{bmatrix} 3 \\ 2 \end{bmatrix}$.
\begin{enumerate}
\setcounter{enumi}{18}
    \item Is $b$ in the set $\{v_1, v_2\}$?
    \item Is $b$ in the span, Span$\{v_1, v_2\}$?
\end{enumerate}

\textbf{Parameterized Span} \\
Find the value(s) of $h$ for which $b$ is in the span of the given vectors.
\begin{enumerate}
\setcounter{enumi}{20}
    \item $v_1 = \begin{bmatrix} 1 \\ -2 \end{bmatrix}, v_2 = \begin{bmatrix} 3 \\ 1 \end{bmatrix}, b = \begin{bmatrix} 5 \\ h \end{bmatrix}$
    \item $v_1 = \begin{bmatrix} 2 \\ -4 \end{bmatrix}, v_2 = \begin{bmatrix} -3 \\ h \end{bmatrix}, b = \begin{bmatrix} 1 \\ -2 \end{bmatrix}$ (For what $h$ is it in the span for ANY $b$?)
    \item $v_1 = \begin{bmatrix} 1 \\ 1 \\ -1 \end{bmatrix}, v_2 = \begin{bmatrix} 2 \\ 3 \\ 0 \end{bmatrix}, b = \begin{bmatrix} 4 \\ 5 \\ h \end{bmatrix}$
    \item $v_1 = \begin{bmatrix} 1 \\ 2 \\ 3 \end{bmatrix}, v_2 = \begin{bmatrix} -1 \\ 0 \\ 1 \end{bmatrix}, v_3 = \begin{bmatrix} 0 \\ 2 \\ h \end{bmatrix}, b = \begin{bmatrix} 2 \\ -4 \\ 6 \end{bmatrix}$. (Find $h$ such that the vectors DO NOT span $\mathbb{R}^3$).
\end{enumerate}

\textbf{Real-World Applications}
\begin{enumerate}
\setcounter{enumi}{24}
    \item \textbf{Center of Mass:} A system consists of three point masses: $2g$ at $\begin{bmatrix} 1 \\ 1 \end{bmatrix}$, $3g$ at $\begin{bmatrix} 5 \\ 0 \end{bmatrix}$, and $5g$ at $\begin{bmatrix} 0 \\ 4 \end{bmatrix}$. Find the coordinates of the center of mass.
    \item \textbf{Center of Mass:} A $4g$ mass is at $v_1 = \begin{bmatrix} 2 \\ 2 \end{bmatrix}$ and a $6g$ mass is at $v_2 = \begin{bmatrix} x \\ y \end{bmatrix}$. If the center of mass of the system is at $\begin{bmatrix} 5 \\ 2 \end{bmatrix}$, find $v_2$.
    \item \textbf{Resource Optimization:} A bakery makes two types of cakes. Cake A requires $2$ cups of flour and $1$ cup of sugar. Cake B requires $3$ cups of flour and $2$ cups of sugar. The bakery has exactly $16$ cups of flour and $9$ cups of sugar. Write a vector equation to find how many of each cake they can bake. 
\end{enumerate}

\textbf{Spanning all of $\mathbb{R}^n$}
\begin{enumerate}
\setcounter{enumi}{27}
    \item Do the vectors $\begin{bmatrix} 1 \\ 0 \end{bmatrix}, \begin{bmatrix} 0 \\ 1 \end{bmatrix}, \begin{bmatrix} -3 \\ 4 \end{bmatrix}$ span all of $\mathbb{R}^2$?
    \item Do the vectors $\begin{bmatrix} 1 \\ 2 \\ 3 \end{bmatrix}$ and $\begin{bmatrix} 4 \\ 5 \\ 6 \end{bmatrix}$ span all of $\mathbb{R}^3$? Why or why not?
\end{enumerate}

\textbf{Advanced Diagnostic Testing: "Find the Flaw"}
\begin{enumerate}
\setcounter{enumi}{29}
    \item \textbf{Flawed logic:} "To check if $b = \begin{bmatrix} 4 \\ 8 \end{bmatrix}$ is in Span$\{v_1, v_2\}$ where $v_1 = \begin{bmatrix} 1 \\ 2 \end{bmatrix}$ and $v_2 = \begin{bmatrix} 2 \\ 4 \end{bmatrix}$, I row reduced the augmented matrix and got $\begin{bmatrix} 1 & 2 & 4 \\ 0 & 0 & 0 \end{bmatrix}$. Because there is a row of zeros, the system is inconsistent, so $b$ is not in the span." Explain the error and provide the correct conclusion.
\end{enumerate}


\newpage
\subsection{Homework 1.4: The Matrix Equation $Ax=b$}

\textbf{Computing $Ax$} \\
Compute the product using the \textbf{Definition} of matrix-vector multiplication (linear combination of columns):
\begin{enumerate}
    \item $\begin{bmatrix} 2 & -3 \\ 1 & 4 \end{bmatrix} \begin{bmatrix} 3 \\ 2 \end{bmatrix}$
    \item $\begin{bmatrix} 1 & 0 & -2 \\ 3 & -1 & 4 \end{bmatrix} \begin{bmatrix} 2 \\ -1 \\ 3 \end{bmatrix}$
    \item $\begin{bmatrix} 4 & -2 \\ 0 & 5 \\ -1 & 1 \end{bmatrix} \begin{bmatrix} -1 \\ 3 \end{bmatrix}$
\end{enumerate}

Compute the product using the \textbf{Row-Vector Rule} (dot product rule):
\begin{enumerate}
\setcounter{enumi}{3}
    \item $\begin{bmatrix} -5 & 2 \\ 1 & -3 \end{bmatrix} \begin{bmatrix} 4 \\ -1 \end{bmatrix}$
    \item $\begin{bmatrix} 1 & 2 & 3 \\ 0 & -1 & 2 \\ 3 & 0 & -4 \end{bmatrix} \begin{bmatrix} -1 \\ 2 \\ 1 \end{bmatrix}$
    \item $\begin{bmatrix} 2 & -1 & 5 \end{bmatrix} \begin{bmatrix} 3 \\ 4 \\ -2 \end{bmatrix}$
\end{enumerate}

\textbf{Dimensional Compatibility} \\
Determine if the product $Ax$ is defined. If undefined, explain why. If defined, state the dimensions of the result.
\begin{enumerate}
\setcounter{enumi}{6}
    \item $A$ is a $4 \times 3$ matrix, $x$ is a $4 \times 1$ vector.
    \item $A = \begin{bmatrix} 1 & 2 \\ 3 & 4 \\ 5 & 6 \end{bmatrix}$, $x = \begin{bmatrix} -1 \\ 1 \end{bmatrix}$.
\end{enumerate}

\textbf{Translation (System/Vector $\leftrightarrow$ Matrix Equation)} \\
Write the given system or vector equation as a matrix equation $Ax=b$.
\begin{enumerate}
\setcounter{enumi}{8}
    \item $5x_1 - 2x_2 = 7$ \\ $-x_1 + 4x_2 = 0$
    \item $x_1 - x_2 + 2x_3 = 4$ \\ $3x_1 + x_3 = -1$ \\ $x_2 - 5x_3 = 2$
    \item $x_1 \begin{bmatrix} 3 \\ -1 \end{bmatrix} + x_2 \begin{bmatrix} 0 \\ 4 \end{bmatrix} + x_3 \begin{bmatrix} -2 \\ 5 \end{bmatrix} = \begin{bmatrix} 8 \\ -3 \end{bmatrix}$
\end{enumerate}

Write the matrix equation as a vector equation, and as a system of linear equations.
\begin{enumerate}
\setcounter{enumi}{11}
    \item $\begin{bmatrix} 1 & 4 \\ -2 & 3 \end{bmatrix} \begin{bmatrix} x_1 \\ x_2 \end{bmatrix} = \begin{bmatrix} 5 \\ 0 \end{bmatrix}$
    \item $\begin{bmatrix} 2 & 0 & -1 \\ 1 & -1 & 4 \\ 0 & 3 & 2 \end{bmatrix} \begin{bmatrix} x \\ y \\ z \end{bmatrix} = \begin{bmatrix} 7 \\ 2 \\ -5 \end{bmatrix}$
    \item $\begin{bmatrix} 1 & 2 & 3 \end{bmatrix} \begin{bmatrix} x_1 \\ x_2 \\ x_3 \end{bmatrix} = \begin{bmatrix} 10 \end{bmatrix}$
\end{enumerate}

\textbf{Solving $Ax=b$} \\
Solve the matrix equation $Ax=b$ by row reducing the augmented matrix. Write the solution as a vector.
\begin{enumerate}
\setcounter{enumi}{14}
    \item $\begin{bmatrix} 1 & -3 \\ 2 & -5 \end{bmatrix} \begin{bmatrix} x_1 \\ x_2 \end{bmatrix} = \begin{bmatrix} 4 \\ 7 \end{bmatrix}$
    \item $\begin{bmatrix} 1 & 2 & -1 \\ 2 & 4 & -2 \end{bmatrix} \begin{bmatrix} x_1 \\ x_2 \\ x_3 \end{bmatrix} = \begin{bmatrix} 3 \\ 6 \end{bmatrix}$
    \item $\begin{bmatrix} 1 & -1 \\ -2 & 2 \\ 3 & -3 \end{bmatrix} \begin{bmatrix} x_1 \\ x_2 \end{bmatrix} = \begin{bmatrix} 2 \\ -4 \\ 5 \end{bmatrix}$
    \item Is $u = \begin{bmatrix} 2 \\ -1 \\ 3 \end{bmatrix}$ in the plane spanned by the columns of $A = \begin{bmatrix} 1 & 0 \\ 0 & 1 \\ 2 & -1 \end{bmatrix}$?
    \item If the augmented matrix $[A \ b]$ reduces to $\begin{bmatrix} 1 & 0 & 4 \\ 0 & 1 & -2 \\ 0 & 0 & 0 \end{bmatrix}$, what is the solution $x$ to $Ax=b$?
\end{enumerate}

\textbf{Existence for all $b$ (Theorem 4)} \\
Determine if the columns of the matrix $A$ span $\mathbb{R}^m$. Justify using Theorem 4.
\begin{enumerate}
\setcounter{enumi}{19}
    \item $A = \begin{bmatrix} 1 & 2 & -1 \\ 0 & 3 & 1 \\ 0 & 0 & 5 \end{bmatrix}$
    \item $A = \begin{bmatrix} 2 & 4 \\ -1 & -2 \\ 3 & 6 \end{bmatrix}$
    \item $A = \begin{bmatrix} 1 & -3 & 2 & 4 \\ -2 & 6 & -4 & -8 \end{bmatrix}$
\end{enumerate}

\textbf{Linearity Properties}
\begin{enumerate}
\setcounter{enumi}{22}
    \item Let $A$ be a $2 \times 2$ matrix. If $A\begin{bmatrix} 1 \\ 0 \end{bmatrix} = \begin{bmatrix} 3 \\ -1 \end{bmatrix}$ and $A\begin{bmatrix} 0 \\ 1 \end{bmatrix} = \begin{bmatrix} -2 \\ 5 \end{bmatrix}$, compute $A\begin{bmatrix} 4 \\ 3 \end{bmatrix}$ using linearity.
    \item If $Ax = b$ and $Ay = 0$, prove that $A(x + 5y) = b$.
    \item If $A(u+v) = \begin{bmatrix} 5 \\ 5 \end{bmatrix}$ and $Au = \begin{bmatrix} 2 \\ 1 \end{bmatrix}$, find $Av$.
\end{enumerate}

\textbf{Parameterized Consistency}
\begin{enumerate}
\setcounter{enumi}{25}
    \item Find all values of $h$ such that $Ax=b$ is consistent: $A = \begin{bmatrix} 1 & -2 \\ -3 & 6 \end{bmatrix}, b = \begin{bmatrix} 4 \\ h \end{bmatrix}$.
    \item Find all values of $k$ such that $A$ spans $\mathbb{R}^2$: $A = \begin{bmatrix} 2 & k \\ 4 & 10 \end{bmatrix}$.
\end{enumerate}

\textbf{Real-World Applications \& Generalizations}
\begin{enumerate}
\setcounter{enumi}{27}
    \item \textbf{Derivative Replication:} In a market with two future states (Up, Down), a risk-free bond pays $\$100$ in both states. A stock pays $\$120$ (Up) and $\$80$ (Down). An option contract pays $\$20$ (Up) and $\$0$ (Down). Write a matrix equation $Ax=b$ to find the number of bonds ($x_1$) and stocks ($x_2$) needed to replicate the option's payoff, and solve it.
    \item \textbf{Generalization:} Explain conceptually why a $3 \times 2$ matrix cannot have columns that span all of $\mathbb{R}^3$.
    \item \textbf{Find the Flaw:} "Since $A$ is a $4 \times 5$ matrix, it has more columns than rows. Therefore, its columns must span $\mathbb{R}^4$." Identify the conceptual error.
\end{enumerate}

\newpage

% ==========================================
% SECTION: DETAILED SOLUTIONS
% ==========================================
\section{Detailed Solutions}

\subsection{Solutions for 1.3: Vector Equations}
\begin{enumerate}
    % 1-3
    \item $u + 2v = \begin{bmatrix} 2 \\ -1 \end{bmatrix} + 2\begin{bmatrix} 3 \\ 4 \end{bmatrix} = \begin{bmatrix} 2+6 \\ -1+8 \end{bmatrix} = \begin{bmatrix} 8 \\ 7 \end{bmatrix}$.
    \item $3u - v + 2w = 3\begin{bmatrix} 2 \\ -1 \end{bmatrix} - \begin{bmatrix} 3 \\ 4 \end{bmatrix} + 2\begin{bmatrix} -1 \\ 2 \end{bmatrix} = \begin{bmatrix} 6-3-2 \\ -3-4+4 \end{bmatrix} = \begin{bmatrix} 1 \\ -3 \end{bmatrix}$.
    \item $c_1\begin{bmatrix} 2 \\ -1 \end{bmatrix} + c_2\begin{bmatrix} -1 \\ 2 \end{bmatrix} = \begin{bmatrix} 7 \\ -5 \end{bmatrix}$. Augmented matrix: $\begin{bmatrix} 2 & -1 & 7 \\ -1 & 2 & -5 \end{bmatrix} \xrightarrow{R_1 \leftrightarrow R_2} \begin{bmatrix} -1 & 2 & -5 \\ 2 & -1 & 7 \end{bmatrix} \xrightarrow{R_2+2R_1} \begin{bmatrix} -1 & 2 & -5 \\ 0 & 3 & -3 \end{bmatrix}$. So $3c_2 = -3 \Rightarrow c_2 = -1$. $-c_1 + 2(-1) = -5 \Rightarrow c_1 = 3$.
    
    % 4-6
    \item $x_1 \begin{bmatrix} 3 \\ 1 \end{bmatrix} + x_2 \begin{bmatrix} -1 \\ 4 \end{bmatrix} = \begin{bmatrix} 5 \\ -2 \end{bmatrix}$.
    \item $x_1 \begin{bmatrix} 1 \\ 2 \\ -1 \end{bmatrix} + x_2 \begin{bmatrix} 2 \\ -1 \\ 0 \end{bmatrix} + x_3 \begin{bmatrix} -1 \\ 3 \\ 5 \end{bmatrix} = \begin{bmatrix} 0 \\ 4 \\ 7 \end{bmatrix}$.
    \item $x_1 \begin{bmatrix} 4 \\ 1 \\ -2 \end{bmatrix} + x_2 \begin{bmatrix} -2 \\ 1 \\ 5 \end{bmatrix} = \begin{bmatrix} 1 \\ 3 \\ 0 \end{bmatrix}$.
    
    % 7-9
    \item $x_1 + 4x_2 = 6$ \\ $-3x_1 + 2x_2 = -1$
    \item $2x_1 - x_2 = 0$ \\ $5x_2 = 2$ \\ $x_1 + 3x_2 = -4$
    \item $x_1 + 4x_3 = 8$ \\ $x_1 - 3x_2 + 2x_3 = 9$ \\ $-2x_1 + x_2 = 10$

    % 10-15
    \item $\begin{bmatrix} 1 & -3 & -5 \\ 2 & 4 & 10 \end{bmatrix} \xrightarrow{R_2-2R_1} \begin{bmatrix} 1 & -3 & -5 \\ 0 & 10 & 20 \end{bmatrix}$. Consistent ($x_2=2, x_1=1$). \textbf{Yes}.
    \item $\begin{bmatrix} 2 & -1 & 3 \\ 4 & -2 & 5 \end{bmatrix} \xrightarrow{R_2-2R_1} \begin{bmatrix} 2 & -1 & 3 \\ 0 & 0 & -1 \end{bmatrix}$. Row of form $[0 \ 0 \ | \ -1]$ is inconsistent. \textbf{No}.
    \item $\begin{bmatrix} 1 & 2 & 4 \\ 0 & 1 & 1 \\ -1 & 3 & 1 \end{bmatrix} \xrightarrow{R_3+R_1} \begin{bmatrix} 1 & 2 & 4 \\ 0 & 1 & 1 \\ 0 & 5 & 5 \end{bmatrix} \xrightarrow{R_3-5R_2} \begin{bmatrix} 1 & 2 & 4 \\ 0 & 1 & 1 \\ 0 & 0 & 0 \end{bmatrix}$. Consistent. \textbf{Yes}.
    \item $\begin{bmatrix} 1 & 3 & 5 \\ -1 & 0 & -2 \\ 2 & 1 & 6 \end{bmatrix} \xrightarrow{R_2+R_1, R_3-2R_1} \begin{bmatrix} 1 & 3 & 5 \\ 0 & 3 & 3 \\ 0 & -5 & -4 \end{bmatrix} \xrightarrow{R_2/3} \begin{bmatrix} 1 & 3 & 5 \\ 0 & 1 & 1 \\ 0 & -5 & -4 \end{bmatrix} \xrightarrow{R_3+5R_2} \begin{bmatrix} 1 & 3 & 5 \\ 0 & 1 & 1 \\ 0 & 0 & 1 \end{bmatrix}$. Inconsistent. \textbf{No}.
    \item $\begin{bmatrix} 1 & 1 & 1 & 4 \\ 0 & 1 & 1 & -2 \\ 0 & 0 & 1 & 7 \end{bmatrix}$. Matrix is already in echelon form, no contradictory rows. \textbf{Yes}.
    \item Notice $a_2 = -a_1$ and $a_3 = 2a_1$. They are all scalar multiples of $\begin{bmatrix} 1 \\ 2 \\ 3 \end{bmatrix}$. Since $b = \begin{bmatrix} 1 \\ 1 \\ 1 \end{bmatrix}$ is not a multiple of $\begin{bmatrix} 1 \\ 2 \\ 3 \end{bmatrix}$, it is not in the span. \textbf{No}.

    % 16-18
    \item Notice $v_1 = -3v_2$. They are scalar multiples (collinear), so their span is a \textbf{line} through the origin in $\mathbb{R}^3$.
    \item Because $v_3 = v_1 + v_2$, $v_3$ is already in the span of $v_1, v_2$. The span is completely determined by the two non-collinear vectors $v_1, v_2$. Geometrically, this forms a \textbf{plane} through the origin.
    \item \textbf{No}. The zero vector can be formed by multiplying any vector by the scalar 0. Including it in a set does not add any new direction or dimension to the linear combinations.

    % 19-20
    \item \textbf{No.} The set $\{v_1, v_2\}$ contains exactly two specific elements. $b$ is neither $v_1$ nor $v_2$.
    \item Check $\begin{bmatrix} 1 & 2 & 3 \\ 3 & -1 & 2 \end{bmatrix} \xrightarrow{R_2-3R_1} \begin{bmatrix} 1 & 2 & 3 \\ 0 & -7 & -7 \end{bmatrix}$. Consistent ($c_2=1, c_1=1$). \textbf{Yes}, it is in the span.

    % 21-24
    \item $\begin{bmatrix} 1 & 3 & 5 \\ -2 & 1 & h \end{bmatrix} \xrightarrow{R_2+2R_1} \begin{bmatrix} 1 & 3 & 5 \\ 0 & 7 & h+10 \end{bmatrix}$. There is no restriction on $h$ that causes a contradiction (no $[0 \ 0 \ | \ c]$). It is in the span for \textbf{all real $h$}.
    \item Notice $v_1 = \begin{bmatrix} 2 \\ -4 \end{bmatrix}$. For the span to cover ALL $\mathbb{R}^2$, $v_1$ and $v_2$ must not be collinear. If $v_2 = c \cdot v_1$, they span a line. Setting $-3 = c(2) \Rightarrow c = -1.5$. So $h = -1.5(-4) = 6$. For the span to be all $\mathbb{R}^2$, \textbf{$h \neq 6$}.
    \item $\begin{bmatrix} 1 & 2 & 4 \\ 1 & 3 & 5 \\ -1 & 0 & h \end{bmatrix} \xrightarrow{R_2-R_1, R_3+R_1} \begin{bmatrix} 1 & 2 & 4 \\ 0 & 1 & 1 \\ 0 & 2 & h+4 \end{bmatrix} \xrightarrow{R_3-2R_2} \begin{bmatrix} 1 & 2 & 4 \\ 0 & 1 & 1 \\ 0 & 0 & h+2 \end{bmatrix}$. For consistency, we need $h+2 = 0 \Rightarrow \textbf{h = -2}$.
    \item They do not span $\mathbb{R}^3$ if they are linearly dependent (i.e., missing a pivot in row 3). $\begin{bmatrix} 1 & -1 & 0 \\ 2 & 0 & 2 \\ 3 & 1 & h \end{bmatrix} \xrightarrow{R_2-2R_1, R_3-3R_1} \begin{bmatrix} 1 & -1 & 0 \\ 0 & 2 & 2 \\ 0 & 4 & h \end{bmatrix} \xrightarrow{R_3-2R_2} \begin{bmatrix} 1 & -1 & 0 \\ 0 & 2 & 2 \\ 0 & 0 & h-4 \end{bmatrix}$. For them NOT to span, row 3 must not have a pivot, so $\textbf{h = 4}$.

    % 25-27
    \item $CM = \frac{1}{2+3+5} \left( 2\begin{bmatrix} 1 \\ 1 \end{bmatrix} + 3\begin{bmatrix} 5 \\ 0 \end{bmatrix} + 5\begin{bmatrix} 0 \\ 4 \end{bmatrix} \right) = \frac{1}{10} \begin{bmatrix} 2+15+0 \\ 2+0+20 \end{bmatrix} = \begin{bmatrix} 17/10 \\ 22/10 \end{bmatrix} = \textbf{\begin{bmatrix} 1.7 \\ 2.2 \end{bmatrix}}$.
    \item $\frac{1}{4+6} \left( 4\begin{bmatrix} 2 \\ 2 \end{bmatrix} + 6\begin{bmatrix} x \\ y \end{bmatrix} \right) = \begin{bmatrix} 5 \\ 2 \end{bmatrix} \Rightarrow \begin{bmatrix} 8 \\ 8 \end{bmatrix} + \begin{bmatrix} 6x \\ 6y \end{bmatrix} = \begin{bmatrix} 50 \\ 20 \end{bmatrix} \Rightarrow \begin{bmatrix} 6x \\ 6y \end{bmatrix} = \begin{bmatrix} 42 \\ 12 \end{bmatrix} \Rightarrow \textbf{v_2 = \begin{bmatrix} 7 \\ 2 \end{bmatrix}}$.
    \item Let $x_1$ be Cake A and $x_2$ be Cake B. Vector eq: $\textbf{x_1 \begin{bmatrix} 2 \\ 1 \end{bmatrix} + x_2 \begin{bmatrix} 3 \\ 2 \end{bmatrix} = \begin{bmatrix} 16 \\ 9 \end{bmatrix}}$.

    % 28-30
    \item \textbf{Yes.} The matrix formed by the first two vectors is the identity matrix $I_2$, which has a pivot in every row. Any additional vectors keep the span as $\mathbb{R}^2$.
    \item \textbf{No.} The matrix is $3 \times 2$. It can have at most 2 pivots. Thus, it cannot have a pivot in all 3 rows. By Theorem 4, it cannot span $\mathbb{R}^3$.
    \item \textbf{Flaw:} The student concluded that a zero row ($\begin{bmatrix} 0 & 0 & 0 \end{bmatrix}$) means inconsistency. This is false. A zero row indicates $0=0$, which is a true statement (free variable / redundant equation). Inconsistency only occurs with $[0 \ 0 \ | \ c]$ where $c \neq 0$. Here, $4=2(2)$, so $b=2v_2$. The system is consistent, and $b$ \textbf{is} in the span.
\end{enumerate}

\subsection{Solutions for 1.4: The Matrix Equation $Ax=b$}
\begin{enumerate}
    % 1-3
    \item $3\begin{bmatrix} 2 \\ 1 \end{bmatrix} + 2\begin{bmatrix} -3 \\ 4 \end{bmatrix} = \begin{bmatrix} 6 \\ 3 \end{bmatrix} + \begin{bmatrix} -6 \\ 8 \end{bmatrix} = \textbf{\begin{bmatrix} 0 \\ 11 \end{bmatrix}}$.
    \item $2\begin{bmatrix} 1 \\ 3 \end{bmatrix} - 1\begin{bmatrix} 0 \\ -1 \end{bmatrix} + 3\begin{bmatrix} -2 \\ 4 \end{bmatrix} = \begin{bmatrix} 2 \\ 6 \end{bmatrix} + \begin{bmatrix} 0 \\ 1 \end{bmatrix} + \begin{bmatrix} -6 \\ 12 \end{bmatrix} = \textbf{\begin{bmatrix} -4 \\ 19 \end{bmatrix}}$.
    \item $-1\begin{bmatrix} 4 \\ 0 \\ -1 \end{bmatrix} + 3\begin{bmatrix} -2 \\ 5 \\ 1 \end{bmatrix} = \textbf{\begin{bmatrix} -10 \\ 15 \\ 4 \end{bmatrix}}$.
    
    % 4-6
    \item $\begin{bmatrix} (-5)(4) + (2)(-1) \\ (1)(4) + (-3)(-1) \end{bmatrix} = \textbf{\begin{bmatrix} -22 \\ 7 \end{bmatrix}}$.
    \item $\begin{bmatrix} 1(-1)+2(2)+3(1) \\ 0(-1)-1(2)+2(1) \\ 3(-1)+0(2)-4(1) \end{bmatrix} = \textbf{\begin{bmatrix} 6 \\ 0 \\ -7 \end{bmatrix}}$.
    \item $[2(3) + (-1)(4) + 5(-2)] = [6 - 4 - 10] = \textbf{\begin{bmatrix} -8 \end{bmatrix}}$.

    % 7-8
    \item \textbf{Undefined.} The number of columns in $A$ (3) does not match the number of entries in $x$ (4).
    \item \textbf{Defined.} $A$ is $3 \times 2$ and $x$ is $2 \times 1$. The result will be a $\mathbf{3 \times 1}$ vector.

    % 9-14
    \item $\textbf{\begin{bmatrix} 5 & -2 \\ -1 & 4 \end{bmatrix} \begin{bmatrix} x_1 \\ x_2 \end{bmatrix} = \begin{bmatrix} 7 \\ 0 \end{bmatrix}}$.
    \item $\textbf{\begin{bmatrix} 1 & -1 & 2 \\ 3 & 0 & 1 \\ 0 & 1 & -5 \end{bmatrix} \begin{bmatrix} x_1 \\ x_2 \\ x_3 \end{bmatrix} = \begin{bmatrix} 4 \\ -1 \\ 2 \end{bmatrix}}$.
    \item $\textbf{\begin{bmatrix} 3 & 0 & -2 \\ -1 & 4 & 5 \end{bmatrix} \begin{bmatrix} x_1 \\ x_2 \\ x_3 \end{bmatrix} = \begin{bmatrix} 8 \\ -3 \end{bmatrix}}$.
    \item Vec: $\textbf{x_1 \begin{bmatrix} 1 \\ -2 \end{bmatrix} + x_2 \begin{bmatrix} 4 \\ 3 \end{bmatrix} = \begin{bmatrix} 5 \\ 0 \end{bmatrix}}$. Sys: $\textbf{x_1 + 4x_2 = 5, \ -2x_1 + 3x_2 = 0}$.
    \item Vec: $\textbf{x \begin{bmatrix} 2 \\ 1 \\ 0 \end{bmatrix} + y \begin{bmatrix} 0 \\ -1 \\ 3 \end{bmatrix} + z \begin{bmatrix} -1 \\ 4 \\ 2 \end{bmatrix} = \begin{bmatrix} 7 \\ 2 \\ -5 \end{bmatrix}}$. Sys: $2x-z=7, \ x-y+4z=2, \ 3y+2z=-5$.
    \item Vec: $\textbf{x_1[1] + x_2[2] + x_3[3] = [10]}$. Sys: $\textbf{x_1 + 2x_2 + 3x_3 = 10}$.

    % 15-19
    \item $\begin{bmatrix} 1 & -3 & 4 \\ 2 & -5 & 7 \end{bmatrix} \xrightarrow{R_2-2R_1} \begin{bmatrix} 1 & -3 & 4 \\ 0 & 1 & -1 \end{bmatrix} \xrightarrow{R_1+3R_2} \begin{bmatrix} 1 & 0 & 1 \\ 0 & 1 & -1 \end{bmatrix}$. Vector: $\textbf{\begin{bmatrix} 1 \\ -1 \end{bmatrix}}$.
    \item $\begin{bmatrix} 1 & 2 & -1 & 3 \\ 2 & 4 & -2 & 6 \end{bmatrix} \xrightarrow{R_2-2R_1} \begin{bmatrix} 1 & 2 & -1 & 3 \\ 0 & 0 & 0 & 0 \end{bmatrix}$. Free vars $x_2, x_3$. Solution: $\textbf{\begin{bmatrix} 3 - 2x_2 + x_3 \\ x_2 \\ x_3 \end{bmatrix}}$.
    \item $\begin{bmatrix} 1 & -1 & 2 \\ -2 & 2 & -4 \\ 3 & -3 & 5 \end{bmatrix} \xrightarrow{R_2+2R_1, R_3-3R_1} \begin{bmatrix} 1 & -1 & 2 \\ 0 & 0 & 0 \\ 0 & 0 & -1 \end{bmatrix}$. Row 3 is $0 = -1$. \textbf{No solution}.
    \item This asks if $Ax = u$ is consistent. $\begin{bmatrix} 1 & 0 & 2 \\ 0 & 1 & -1 \\ 2 & -1 & 3 \end{bmatrix} \xrightarrow{R_3-2R_1} \begin{bmatrix} 1 & 0 & 2 \\ 0 & 1 & -1 \\ 0 & -1 & -1 \end{bmatrix} \xrightarrow{R_3+R_2} \begin{bmatrix} 1 & 0 & 2 \\ 0 & 1 & -1 \\ 0 & 0 & -2 \end{bmatrix}$. Inconsistent. \textbf{No.}
    \item The system translates directly to $1x_1 + 0x_2 = 4 \Rightarrow x_1 = 4$, and $0x_1 + 1x_2 = -2 \Rightarrow x_2 = -2$. The solution is $\textbf{\begin{bmatrix} 4 \\ -2 \end{bmatrix}}$.

    % 20-22
    \item $A$ is in echelon form and has pivots in row 1 (1), row 2 (3), row 3 (5). Since every row has a pivot, \textbf{Yes}, it spans $\mathbb{R}^3$.
    \item $\begin{bmatrix} 2 & 4 \\ -1 & -2 \\ 3 & 6 \end{bmatrix} \xrightarrow{R_1 \leftrightarrow R_2} \begin{bmatrix} -1 & -2 \\ 2 & 4 \\ 3 & 6 \end{bmatrix} \xrightarrow{\text{clear col 1}} \begin{bmatrix} -1 & -2 \\ 0 & 0 \\ 0 & 0 \end{bmatrix}$. Missing pivots in rows 2 and 3. \textbf{No}.
    \item $\begin{bmatrix} 1 & -3 & 2 & 4 \\ -2 & 6 & -4 & -8 \end{bmatrix} \xrightarrow{R_2+2R_1} \begin{bmatrix} 1 & -3 & 2 & 4 \\ 0 & 0 & 0 & 0 \end{bmatrix}$. Missing a pivot in row 2. \textbf{No}.

    % 23-25
    \item $A\begin{bmatrix} 4 \\ 3 \end{bmatrix} = 4A\begin{bmatrix} 1 \\ 0 \end{bmatrix} + 3A\begin{bmatrix} 0 \\ 1 \end{bmatrix} = 4\begin{bmatrix} 3 \\ -1 \end{bmatrix} + 3\begin{bmatrix} -2 \\ 5 \end{bmatrix} = \begin{bmatrix} 12 - 6 \\ -4 + 15 \end{bmatrix} = \textbf{\begin{bmatrix} 6 \\ 11 \end{bmatrix}}$.
    \item By linearity, $A(x+5y) = Ax + 5(Ay)$. Substituting the givens: $b + 5(0) = \textbf{b}$. 
    \item $A(u+v) = Au + Av \Rightarrow \begin{bmatrix} 5 \\ 5 \end{bmatrix} = \begin{bmatrix} 2 \\ 1 \end{bmatrix} + Av \Rightarrow Av = \begin{bmatrix} 5 \\ 5 \end{bmatrix} - \begin{bmatrix} 2 \\ 1 \end{bmatrix} = \textbf{\begin{bmatrix} 3 \\ 4 \end{bmatrix}}$.

    % 26-27
    \item $\begin{bmatrix} 1 & -2 & 4 \\ -3 & 6 & h \end{bmatrix} \xrightarrow{R_2+3R_1} \begin{bmatrix} 1 & -2 & 4 \\ 0 & 0 & h+12 \end{bmatrix}$. For consistency, $h+12 = 0 \Rightarrow \textbf{h = -12}$.
    \item To span $\mathbb{R}^2$, the matrix must have a pivot in every row, meaning the columns cannot be scalar multiples of each other. $2 \times 2 = 4$, so we need $k \times 2 \neq 10 \Rightarrow \textbf{k \neq 5}$.

    % 28-30
    \item Matrix $A = \begin{bmatrix} 100 & 120 \\ 100 & 80 \end{bmatrix}$, $b = \begin{bmatrix} 20 \\ 0 \end{bmatrix}$. $\begin{bmatrix} 100 & 120 & 20 \\ 100 & 80 & 0 \end{bmatrix} \xrightarrow{R_2-R_1} \begin{bmatrix} 100 & 120 & 20 \\ 0 & -40 & -20 \end{bmatrix}$. $x_2 = 1/2$. $100x_1 + 120(0.5) = 20 \Rightarrow 100x_1 = -40 \Rightarrow x_1 = -0.4$. Portfolio: \textbf{Buy 0.5 shares of stock, borrow 0.4 bonds.}
    \item A $3 \times 2$ matrix has 2 columns. Since each column can supply at most 1 pivot, the matrix can have at most \textbf{2 pivots}. To span $\mathbb{R}^3$, Theorem 4 requires a pivot in \textbf{every row} (3 pivots). Because $2 < 3$, it can never span $\mathbb{R}^3$.
    \item \textbf{Flaw:} Having more columns than rows ($n > m$) guarantees free variables (meaning if a solution exists, there are infinitely many), but it does \textbf{not} guarantee a pivot in every row. For example, a $4 \times 5$ matrix could have a row of zeros. Spanning requires $m$ pivots, which depends on the actual entries, not just the dimensions.
\end{enumerate}

\newpage

% ==========================================
% APPENDIX: CONCEPT CHECKLIST \& MAPPING
% ==========================================
\section{Appendix: Concept Checklist \& Question Mapping}

\subsection{Section 1.3: Vector Equations}
\begin{itemize}
    \item \textbf{Vector Arithmetic \& Linear Combinations:} 1.3.1, 1.3.2, 1.3.3
    \item \textbf{Translation (System $\leftrightarrow$ Vector Equation):} 1.3.4, 1.3.5, 1.3.6, 1.3.7, 1.3.8, 1.3.9
    \item \textbf{Membership in a Span:} 1.3.10, 1.3.11, 1.3.12, 1.3.13, 1.3.14, 1.3.15
    \item \textbf{Geometric Interpretation:} 1.3.16, 1.3.17, 1.3.18
    \item \textbf{Set vs. Span:} 1.3.19, 1.3.20
    \item \textbf{Parameterized Span:} 1.3.21, 1.3.22, 1.3.23, 1.3.24
    \item \textbf{Real-World Applications:} 1.3.25, 1.3.26, 1.3.27
    \item \textbf{Spanning all of $\mathbb{R}^n$:} 1.3.28, 1.3.29
    \item \textbf{Advanced Diagnostic ("Find the Flaw"):} 1.3.30
\end{itemize}

\subsection{Section 1.4: The Matrix Equation $Ax=b$}
\begin{itemize}
    \item \textbf{Computing $Ax$ (Definition \& Row-Vector Rule):} 1.4.1, 1.4.2, 1.4.3, 1.4.4, 1.4.5, 1.4.6
    \item \textbf{Dimensional Compatibility:} 1.4.7, 1.4.8
    \item \textbf{Translation (System/Vector $\leftrightarrow$ Matrix Equation):} 1.4.9, 1.4.10, 1.4.11, 1.4.12, 1.4.13, 1.4.14
    \item \textbf{Solving $Ax=b$:} 1.4.15, 1.4.16, 1.4.17, 1.4.18, 1.4.19
    \item \textbf{Existence for all $b$ (Theorem 4):} 1.4.20, 1.4.21, 1.4.22
    \item \textbf{Linearity Properties:} 1.4.23, 1.4.24, 1.4.25
    \item \textbf{Parameterized Consistency:} 1.4.26, 1.4.27
    \item \textbf{Real-World Applications:} 1.4.28
    \item \textbf{Generalization \& Diagnostics:} 1.4.29, 1.4.30
\end{itemize}

\end{document}